\documentclass[11pt]{article}
\usepackage[utf8]{inputenc}
\usepackage{amsmath,amssymb}
\usepackage{geometry}
\usepackage{graphicx}
\geometry{margin=1in}

\title{Why Element-wise Array MSE Collapses to Low-Variance Prediction}
\author{}
\date{}

\begin{document}
\maketitle

\paragraph{一句话结论}
关键原因就是：你在做\textbf{数组每个元素的 MSE 回归}，其统计最优解是逐维条件均值，条件均值变化弱时就会收缩成低方差窄带。

\paragraph{(1) 你的训练目标}
设 $Y=(Y_1,\dots,Y_K)$ 为 $K$ 维 regret 向量，预测为 $f(X)=(f_1(X),\dots,f_K(X))$。你优化的是
\[
\mathcal L(f)=\mathbb E\|Y-f(X)\|_2^2
=\sum_{k=1}^K\mathbb E\big[(Y_k-f_k(X))^2\big].
\]
这等价于把每个维度 $k$ 当作独立标量回归。

\paragraph{(2) 每一维 MSE 的解析最优解}
对每个 $k$，总体风险最优为
\[
f_k^*(X)=\mathbb E[Y_k\mid X].
\]
所以整体最优是
\[
f^*(X)=\mathbb E[Y\mid X].
\]
若特征对目标的信息弱，则 $\mathbb E[Y_k\mid X]\approx c_k$（近常数），故预测集中在窄区间。

\paragraph{(3) 为什么是低方差（方差分解）}
每一维满足
\[
\mathrm{Var}(Y_k)
=\mathbb E[\mathrm{Var}(Y_k\mid X)]
+\mathrm{Var}(\mathbb E[Y_k\mid X]).
\]
模型最多只能学到第二项。若第二项小，则最优预测本身方差就小，出现你看到的“横向一条带”。

\paragraph{图像证据}
\begin{figure}[h!]
    \centering
    \includegraphics[width=0.82\linewidth]{../Feb20_appnp_m1_regret/scatter_all_pairs.png}
    \caption{All pairs true vs predicted: predicted values are compressed into a narrow band.}
\end{figure}

\begin{figure}[h!]
    \centering
    \includegraphics[width=0.85\linewidth]{opf_vA_balance_test_adjmode_real_true_regret_heatmap_sorted.png}
    \caption{True regret matrix (all data): target has wide/heterogeneous structure.}
\end{figure}

\paragraph{补充}
归一化会改善优化数值条件，但不是这次塌缩的主因；主因是 ``逐元素 MSE $\Rightarrow$ 条件均值解'' 在当前特征信息下天然低方差。

\end{document}
